\chapter{LİTERATÜR TARAMASI}
\label{ch:literatur}

Bu bölümde çalışmanın dayandığı dört alan özetlenmiştir: destek talebi sınıflandırma, erişim
destekli üretim (RAG), büyük dil modellerinin güvenilirliği ve güvenliği ile küçük örneklerle dürüst
ölçme.

\section{Destek Talebi Sınıflandırma}

Yardım masası taleplerinin otomatik sınıflandırılması uzun süredir klasik makine öğrenmesi
yöntemleriyle (TF-IDF, destek vektör makineleri, derin sinir ağları) çalışılmaktadır. Revina ve
arkadaşları \cite{revina2020}, BT destek talepleri üzerinde yaptıkları karşılaştırmada basit ve
yorumlanabilir yöntemlerin çoğu zaman karmaşık modeller kadar iyi sonuç verdiğini göstermiştir. Bu
yöntemlerin ortak gereksinimi, binlerce etiketli örnektir. Büyük dil modelleri ise az sayıda örnekle
veya yalnızca bir talimat metniyle (\textit{zero-shot}) sınıflandırma yapabilmektedir \cite{brown2020}.
BŞEÜ için etiketli geçmiş talep verisi bulunmadığından bu çalışmada ikinci yol seçilmiş, etiket kural
kitabı doğrudan modelin talimatı (system prompt) olarak kullanılmıştır. Modelin etiketten önce
gerekçesini yazması, düşünce zinciri (\textit{chain-of-thought}) yaklaşımının \cite{wei2022} basit bir
uygulamasıdır.

\section{Erişim Destekli Üretim (RAG)}

Lewis ve arkadaşları \cite{lewis2020}, bir dil modelinin cevabını harici bir belge deposundan
getirilen parçalara dayandırmayı (RAG) önermiştir. Böylece modelin bilmediği veya güncel olmayan
bilgi yerine, kurumun kendi belgeleri kullanılır ve cevap kaynağıyla birlikte gösterilebilir. Belge
arama genellikle yoğun vektör temsilleriyle (\textit{dense retrieval}) yapılır \cite{karpukhin2020};
çok büyük koleksiyonlarda FAISS gibi yaklaşık en yakın komşu kütüphaneleri kullanılır
\cite{johnson2019}. Gao ve arkadaşlarının derlemesi \cite{gao2023}, anahtar kelime ve vektör aramanın
birleştirildiği karma (hibrit) aramanın ve yeniden sıralamanın (\textit{reranking}) doğruluğu
artırdığını belirtmektedir. İki sıralı listeyi birleştirmek için yaygın ve parametresiz bir yöntem
karşılıklı sıra birleştirmedir (RRF) \cite{cormack2009}. RAG sistemlerinin değerlendirilmesi için
RAGAS gibi otomatik ölçütler önerilmiştir \cite{es2024}.

\section{Dil Modellerinin Güvenilirliği ve Güvenliği}

Ji ve arkadaşlarının derlemesi \cite{ji2023}, dil modellerinin kaynakta bulunmayan bilgiyi akıcı
biçimde üretebildiğini (halüsinasyon) ve bunun özellikle bilgi gerektiren görevlerde ciddi bir sorun
olduğunu göstermektedir. Kullanıcı metninin modele girdi olarak verildiği uygulamalarda ise
\textit{prompt injection} saldırıları ortaya çıkar: Perez ve Ribeiro \cite{perez2022} talimatların
``önceki talimatları unut'' gibi ifadelerle ele geçirilebildiğini, Greshake ve arkadaşları
\cite{greshake2023} ise bu saldırının dolaylı yollardan da yapılabildiğini göstermiştir. OWASP, bu
riski büyük dil modeli uygulamaları için en önemli risk olarak listelemektedir \cite{owasp2025}. Bu
nedenle bu çalışmada modelin çıktısına güvenmek yerine, çıktı kodla denetlenmiştir: alıntılanan kural
belgede birebir yoksa cevap gönderilmez.

Ücretsiz katmandaki model servisleri sık sık yoğunluk nedeniyle hata verir veya yavaş cevap verir.
Dean ve Barroso \cite{dean2013}, büyük ölçekli sistemlerde yavaş isteklerin (\textit{tail latency})
etkisini azaltmak için belirli bir süre sonra aynı isteği ikinci bir kaynağa da gönderme
(\textit{hedged request}) yöntemini önermiştir. Canlı sistemdeki model zinciri bu fikre dayanır.

\section{Küçük Örneklerle Ölçme}

Değerlendirme setleri küçük olduğunda tek bir yüzde değeri yanıltıcıdır. Wilson \cite{wilson1927}
güven aralığı, özellikle oran 0 veya 1'e yakınken normal yaklaşıma göre daha doğru aralık verir ve bu
çalışmadaki bütün ölçütlerde kullanılmıştır. Ayrıca istem (prompt) geliştirmede kullanılan verinin
skoru iyimser göstereceği bilindiğinden, ayar için hiç kullanılmayan kilitli test setleri
ayrılmıştır. Cevap kalitesinin bir başka dil modeliyle puanlanması (\textit{LLM-as-a-judge})
\cite{zheng2023} kalan dönem için planlanmıştır.

\section{Literatür Özeti}

Tablo~\ref{tab:literatur} incelenen çalışmaları ve bu projeye katkılarını özetlemektedir.

\begin{table}[!ht]
\centering
\caption{Literatür özeti ve bu çalışmaya yansıması}
\label{tab:literatur}
\footnotesize
\renewcommand{\arraystretch}{1.25}
\begin{tabularx}{\textwidth}{@{}p{2.5cm} p{2.6cm} X X@{}}
\toprule
\textbf{Yazar (Yıl)} & \textbf{Konu} & \textbf{Bulgu} & \textbf{Bu çalışmada} \\
\midrule
Lewis vd. (2020) \cite{lewis2020} & RAG & Cevabı getirilen belgelere dayandırmak doğruluğu artırır & Cevaplar 28 resmî belgeden üretilir \\
Karpukhin vd. (2020) \cite{karpukhin2020} & Yoğun vektör arama & Anlam vektörleri anahtar kelimeyi geçer & 768 boyutlu çok dilli gömme \\
Cormack vd. (2009) \cite{cormack2009} & RRF & Basit sıra birleştirme güçlü sonuç verir & Karma arama için planlandı \\
Revina vd. (2020) \cite{revina2020} & Talep sınıflandırma & Basit yöntemler çoğu zaman yeterli & Az sayıda kural, açık kod \\
Wei vd. (2022) \cite{wei2022} & Düşünce zinciri & Önce gerekçe, sonra cevap & Model önce gerekçe/alıntı yazar \\
Ji vd. (2023) \cite{ji2023} & Halüsinasyon & Model kaynaksız bilgi uydurabilir & Birebir alıntı kodla denetlenir \\
Greshake vd. (2023) \cite{greshake2023} & Prompt injection & Girdi metni modeli yönlendirebilir & Talep çitlenir, 12 saldırı testi \\
Dean, Barroso (2013) \cite{dean2013} & Kuyruk gecikmesi & Yavaş isteğe ikinci kopya gönderilir & 9 sn sonra ikinci modelle yarış \\
Wilson (1927) \cite{wilson1927} & Güven aralığı & Küçük örnekte doğru aralık & Her ölçüt \%95 GA ile \\
\bottomrule
\end{tabularx}
\end{table}
